% !TeX root = ../main.tex

\chapter{Conclusion and Future Work}\label{chapter:conclusion}

\section{Summary}

This thesis introduced InsertGS, a three-stage pipeline for inserting
mesh-based objects into 3D Gaussian Splatting reconstructions of indoor
scenes. By extracting an IBL cubemap from the host 3DGS scene with proper
EWA splatting, rendering the inserted object under that environment in
Blender, and depth-gated alpha compositing into novel views, the system
produces insertions whose shading and contact shadows are consistent with
the captured room, while leaving the host scene unchanged and compatible
with existing real-time 3DGS viewers.

\section{Limitations}

\begin{itemize}
  \item \textbf{Static illumination.} The cubemap is captured once per
        insertion point; moving the object across a room with strongly
        varying lighting requires re-baking.
  \item \textbf{Geometry from Gaussians.} The host scene exposes no
        explicit surfaces, so contact shadows rely on a manually placed
        ground plane.
  \item \textbf{Single bounce.} The inserted object does not reflect light
        back onto the room; this would require a differentiable round-trip
        through the Gaussian representation.
\end{itemize}

\section{Future work}

A natural extension is to learn a lightweight surface proxy alongside the
3DGS scene so that the ground plane (and other receivers) can be inferred
automatically. A second direction is to back-propagate the inserted
object's contribution into a small set of Gaussians around it, enabling
mutual illumination without re-training the entire host scene.
